%\pagestyle{mystyle}
\chapter{Entity Extraction Assessment for the Legal Domain}
\label{chap:extraction}
\todo{information extraction or enitty extraction?}
% All the papers on knowledge extraction
% \section{IJCKG and NIL Prediction}
% \section{Knowledge Extraction in Specific Domains: Specific Issues}
% \subsection{Legal Judgments}
% \subsection{Instant Messaging Data in Criminal Investigations}
% \section{Organizing Knowledge in Specific Domains}
% Especially infrastructure to organize and hold data.
% \subsection{Judgments}
% \subsection{IM Criminal Investigations}
% \section{Adapting to Specific Domains}

\todo{present mammotab as additional contribution disjoint from thesis.}
This chapter presents the contributions about \emph{incremental \acl{nel}} (see \Cref{sec:incrementalel}), co-authored during my \phd  These have been published and presented in 2022 at the 11th International Joint Conference on Knowledge Graphs (IJCKG 2022)---\textcite{ijckg2022}---and at the Fourth edition of the Semantic Web Challenge on Tabular Data to Knowledge Graph Matching (SemTab 2022)---\textcite{mammotab}. Both the events were co-located with the 21st International Semantic Web Conference (ISWC 2022).

In 2022, as described in \Cref{sec:incrementalel}, the advancements in the NIL prediction task were left behind the one in \ac{nel}. For instance, the BLINK approach~\cite{blink} was not yet adapted for NIL prediction.
For this reason, we worked towards producing datasets and evaluation procedures for aiding the research in this field, focusing on unstructured text~\cite{ijckg2022} (\Cref{sec:ijckg}), but also on structured tabular data~\cite{mammotab} (\Cref{sec:mammotab})---for which a brief backround review is given in \Cref{sec:sti}.

% \todo{add brief intro with contributions}
% \todo{focus on incremental/nil prediction}

% % abstract from ijckg2022
% In this paper, we
% advocate for batch-based incremental entity extraction methods
% that can exploit NEL with a background KB, detect mentions of
% entities that are not present in the KB yet (NIL mentions), and up-
% date the KB with the novel entities. Based on this assumption, we
% present a methodology to evaluate NEL-based incremental entity
% extraction, which can be applied to a “static” dataset for evaluating
% NEL into a dataset for evaluating incremental entity extraction. We
% apply this methodology to an existing benchmark for evaluating
% NEL algorithms, and evaluate an incremental extraction pipeline
% that orchestrates different strong state-of-the-art and baseline al-
% gorithms for the tasks involved in the extraction process, namely,
% NEL, NIL prediction, and NIL clustering. In presenting our experi-
% ments, we demonstrate the increased difficulty of the information
% extraction task in incremental settings and discuss the strengths of
% the available solutions as well as open challenges.

% % abstract from mammotab
% In this paper, we present MammoTab, a dataset composed of 1M Wikipedia tables extracted from over
% 20M Wikipedia pages and annotated through Wikidata. The lack of this kind of dataset in the state-
% of-the-art makes MammoTab a good resource for testing and training Semantic Table Interpretation
% approaches. The dataset has been designed to cover several key challenges, such as disambiguation,
% homonym and NIL-mentions. The dataset has been evaluated using MTab, one of the best approaches of
% the SemTab challenge.

In particular, in \textcite{ijckg2022}, we focus on batch-based incremental entity extraction, proposing a procedure to adapt existing \ac{nel} datasets to this setting. We apply this method to the WNUM \ac{el} dataset~\cite{Onoe_Durrett_2020} producing its incremental version, \emph{Incremental WNUM}, which is available from GitHub~\footnote{https://github.com/rpo19/Incremental-Entity-Extraction}. This approach leverages \ac{el} with a background \ac{kb}, while also detecting and clustering NIL mentions, and incrementally updating the \ac{kb} with novel entities. We apply this methodology to an existing benchmark, propose a baseline pipeline, and demonstrate the increased difficulty of \acl{ee} in incremental settings.

In \textcite{mammotab}, we introduce \emph{MammoTab}, a large-scale dataset of 1M annotated Wikipedia tables extracted from 20M pages and aligned with Wikidata. This resource was designed to address core challenges in \emph{\acf{sti}}---i.e., the task of identifying entities, classes columns types, and relations between columns in tabular data (see \Cref{sec:sti})---such as entity disambiguation, homonym resolution, and NIL detection.
%Its utility is showcased by evaluating it with one of the leading systems from the SemTab challenge.

Finally, the contributions for incremental \ac{nel} can be summarized as follows:
\begin{enumerate}
    \item A general methodology to adapt any \ac{nel} dataset to incremental \ac{nel}.
    \item A dataset for incremental \ac{el}, suitable for evaluating \ac{nel}, NIL prediction, and NIL clustering---created from an existing \ac{nel} benchmark~\cite{Onoe_Durrett_2020}, applying our methodology.
    \item A baseline pipeline, based on the BLINK bi-encoder~\cite{blink}, which additionally supports NIL prediction and clustering.
    \item Evaluation procedures for incremental \ac{nel}.
    \item An evaluation of the baseline pipeline on the incremental scenario and a discussion of the main challenges.
\end{enumerate}
Similarly the contributions for \ac{sti} are
\begin{itemize}
    \item mammotab with NIL annotations.
    \item evaluation\todo{}
\end{itemize}


%\pagestyle{mystyle}
% \section{Information Extraction in the Legal Domain}
% \chaptermark{KE in Legal}
% \label{chap:legalie}
\todo{kg because in wise it is not only entities but also relations from chat metadata}
\todo{add dave}
\todo{add the end speak about architecture, scalability, adaptability,...}
% \section{User interfaces}
% \subsection{Document Explorer}
% \subsection{Knowledge Editing to Correct Machine Errors}
% Human-in-the-Loop
% \subsection{Faceted Search}
% \subsection{Question Answering with Retrieval Augmented Generation}

% \section{Question Answering with ReFactX}

As anticipated in \Cref{sec:background:ai4legal}, \ac{ai} systems can be applied to the legal domain for addressing the analysis of vast amount of documents, that can come from heterogeneous sources.
During my \phd I worked with two types of documents from the Italian legal domain: civil judgements and investigative chat logs. In both cases, the main problem addressed is the difficulty that the users---such as judges or prosecutors---face in analyzing the high number of records. Indeed, my co-authored contributions seek to empower final users with tools that allow them to efficiently explore and retrieve relevant documents or chat messages.
These proposals have been published in Computer Law \& Security Review (Vol. 52)~\cite{BELLANDI2024}, in the proceedings of the 22nd International Conference of the Italian Association for Artificial Intelligence (AIxIA 2023 -- Advances in Artificial Intelligence)~\cite{aixia_2023}, and of the 25th International Conference on Web Information Systems Engineering (WISE 2024)~\cite{pozzi_wise_2025}.

Data in this domain, as described in \Cref{sec:background:ai4legal}, often contain personal data, and privacy regulations imply important challenges such as the difficulty in obtaining labeled data for evaluating or training \ac{ee} systems. Furthermore, the language used in documents---such as civil judgments---differs from daily language~\cite{ccetindaug2023name_legal_turkish}, while open datasets used for training neural models, such as Common Crawl~\cite{commoncrawl} or The Pile~\cite{gao2020pile800gbdatasetdiverse}, are generally obtained from Web data. Anovil judgments, due to the lack of labeled domain data, a benchmark dataset has been annotated with a semi-automatic approach and a \ac{nel} pipeline, inspired from the one introduced in \Cref{sec:incrementalel:pipeline}, has been evaluated on the labeled data~\cite{BELLANDI2024}. Additionally, different adaptations techniques have been studied to improve \ac{ner} performance on legal data~\cite{aixia_2023}, discussing both the accuracy and the required computation.

For criminal investigation data, we focused the extraction and organization of the entities from instant messaging data. Messages---including voice transcriptions---have been modeled in a knowledge graph that capture the network of contacts of the main suspect. Furthermore, an improved version of the \ac{nel} pipeline from \Cref{sec:incrementalel:pipeline} have been evaluated on data---previously labeled with a similar procedure to the one applied to legal judgments~\cite{pozzi_wise_2025}.

These works also include contributions in \acl{ia}, which are build upon the knowledge extracted from the domain data, and are described later in the next \Cref{chap:legalia}.

In summary, the main contributions in legal \ac{ee} of this \phd are the following.
\begin{itemize}
    \item Creation of annotated benchmark datasets for civil judgments and investigative chat data, enabling structured evaluation of entity extraction methods.
    \item Development and evaluation of \ac{ner} and incremental \ac{nel} pipelines tailored for legal and investigative domains, including a study on domain adaptation techniques.
    \item Knowledge organization using knowledge graphs to represent entities and relationships from investigative chat records.
\end{itemize}

Another consequent issue is the impossibility of using APIs, such as OpenAI ones\footnote{https://openai.com/api/}\todo{mention apis, openai,... in llms}, with legal data that contain personal information---unless previously anonymized or pseudonymized. Moreover, using these APIs may involve transferring sensitive data outside the European Union, raising additional compliance concerns~\cite[Article~44]{gdpr2016}. As a result, institutions are often forced to rely on local models, which, while enabling full control over data storage and security in line with \textcite[Article~32]{gdpr2016}, can be limited by the local infrastructure.

In the remaining of this chapter, I present the contributions of my \phd work on entity extraction in the Italian legal domain.
Before describing these contributions, a review of the related literature on legal \ac{ie} and on \ac{ie} in Italian is provided, in order to contextualize the research.

With respect to civil judgments, due to the lack of labeled domain data, a benchmark dataset has been annotated with a semi-automatic approach and a \ac{nel} pipeline, inspired from the one introduced in \Cref{sec:incrementalel:pipeline}, has been evaluated on the labeled data~\cite{BELLANDI2024}. Additionally, different adaptations techniques have been studied to improve \ac{ner} performance on legal data~\cite{aixia_2023}, discussing both the accuracy and the required computation.

For criminal investigation data, we focused the extraction and organization of the entities from instant messaging data. Messages---including voice transcriptions---have been modeled in a knowledge graph that capture the network of contacts of the main suspect. Furthermore, an improved version of the \ac{nel} pipeline from \Cref{sec:incrementalel:pipeline} have been evaluated on data---previously labeled with a similar procedure to the one applied to legal judgments~\cite{pozzi_wise_2025}.

These works also include contributions in \acl{ia}, which are build upon the knowledge extracted from the domain data, and are described later in the next \Cref{chap:legalia}.

In summary, the main contributions in legal \ac{ee} of this \phd are the following.
\begin{itemize}
    \item Creation of annotated benchmark datasets for civil judgments and investigative chat data, enabling structured evaluation of entity extraction methods.
    \item Development and evaluation of \ac{ner} and incremental \ac{nel} pipelines tailored for legal and investigative domains, including a study on domain adaptation techniques.
    \item Knowledge organization using knowledge graphs to represent entities and relationships from investigative chat records.
\end{itemize}